%!TEX root = ../../main.tex
%% ===============================================================
%% 6.6 플레이어 상호작용 그래프
%% 파일: chapters/06_features/06_interaction_graph.tex
%% 상위: chapters/06_features/chapter.tex
%% 정의 라벨: sec:feat-graph, eq:adjacency
%% 참조 라벨: -
%% 포함 객체: -
%% 인용: -
%% 상태: 초안 (docs/OUTLINE.md 참고)
%% ===============================================================

\section{플레이어 상호작용 그래프}
\label{sec:feat-graph}

구간 $\ell$마다 열 명의 참가자를 정점으로 하는 완전 가중 그래프 $G_\ell = (V, A_\ell)$을 만든다. 인접 가중치는 가우스 커널
\begin{equation}
  A_{ij}(\ell) = \exp\!\left(-\frac{\|\mathrm{pos}_i(\ell) - \mathrm{pos}_j(\ell)\|_2^2}{2\sigma^2}\right)
  \label{eq:adjacency}
\end{equation}
로 정의하며, $\sigma$는 고정값 또는 구간 내 평균 거리 $\bar d$에 비례하는 적응값($\sigma = 0.5\,\bar d$, GraphSAGE 실험)을 쓴다. 사망한 참가자에 연결된 변은 0으로 마스킹하고 자기 루프는 유지한다. MPNN을 위한 변 특징은 $e_{ij} = [\Delta x, \Delta y, d, \log(1+A_{ij})] \in \R^4$이다. 역할 인식 처치(T5)는 $A'_{ij} = A_{ij}\cdot\softplus(R_{\mathrm{role}(i),\mathrm{role}(j)})$로 학습 가능한 $5\times 5$ 역할 상호작용 행렬 $R$을 곱하며, 다중 스케일 변형은 여러 $\sigma_k$의 커널을 학습 가중치로 합한다.
